\bookchapter{Bab 8: Komunikasi: SMS, Telepon, dan Media Sosial}

\chapterintro{Komunikasi adalah misi paling awal dan paling inti dari ponsel. Kini, komunikasi sudah berevolusi dari sekadar telepon dan SMS menjadi ekosistem kompleks yang mencakup pesan instan, media sosial, dan email.}

\booksection{Fitur Lanjutan App Telepon}

App telepon di ponsel modern sudah jauh melampaui sekadar papan tombol. App ini menggabungkan alat-alat canggih seperti perekaman panggilan, pemblokiran gangguan, dan identifikasi pintar.

\booksubsection{Perekaman panggilan (perhatikan kepatuhan hukum):}

\begin{itemize}
  \item Fungsi: merekam isi percakapan penting untuk notulen kerja, wawancara, janji klien, atau penyimpanan bukti sengketa.
  \item Catatan hukum: Di Tiongkok, rekaman panggilan dapat dipakai sebagai bukti bila mendapat persetujuan pihak lain, atau digunakan untuk melindungi hak dan kepentingan sah diri sendiri.
  \item Rekaman sepihak vs persetujuan dua pihak: Secara hukum, rekaman yang diperoleh "tanpa persetujuan pihak lain" tidak otomatis batal. Menurut penafsiran yudisial Mahkamah Agung, selama cara memperoleh rekaman tidak melanggar hak sah orang lain secara serius, tidak melanggar ketentuan hukum yang melarang, dan tidak bertentangan dengan ketertiban umum, rekaman itu masih dapat diajukan sebagai bukti. Namun, sebagian pengadilan provinsi menerapkan standar yang lebih ketat terhadap "rekaman diam-diam" dalam perkara perdata; perkara pidana biasanya memerlukan tinjauan legalitas yang lebih ketat.
  \item Perbedaan antardaerah: Pengadilan di Beijing, Shanghai, Guangdong, dan sekitarnya lebih lunak terhadap rekaman sepihak sebagai bukti perdata; sementara rekaman yang melibatkan rahasia dagang atau privasi pribadi perlu sangat berhati-hati di provinsi mana pun, karena bisa digugurkan karena melanggar hak privasi.
  \item Cara melakukannya: Android biasanya punya tombol "Rekam" yang jelas di layar panggilan; iOS tidak menyediakan fitur perekaman panggilan bawaan, jadi perlu App pihak ketiga atau headset Bluetooth dengan fitur rekam.
\end{itemize}

\begin{tcbwarn}{Garis merah penggunaan rekaman}
\begin{itemize}
  \item Jangan gunakan rekaman untuk pemerasan, penyebaran, atau penghinaan terbuka terhadap orang lain.
  \item Konten yang melibatkan privasi pihak ketiga (misalnya tidak sengaja merekam percakapan orang lain) tidak boleh disebarluaskan sembarangan.
  \item Untuk situasi penting, sebaiknya beri tahu secara lisan di awal panggilan bahwa "panggilan ini akan direkam"—lebih patuh hukum sekaligus lebih kuat sebagai bukti.
\end{itemize}

\end{tcbwarn}

\booksubsection{Pembahasan mendalam penanganan panggilan gangguan:}

\begin{itemize}
  \item Prinsipnya: App telepon menelusuri basis data nomor yang besar melalui internet, lalu menampilkan "panggilan gangguan", "panggilan penipuan", atau nama perusahaan saat ada panggilan masuk.
  \item Cara mengaktifkan: Di Android, buka App Telepon > Pengaturan > Pemblokiran Gangguan; di iOS, buka Pengaturan > Telepon > Pemblokiran \& Identifikasi Pemanggil. Jalur pengaturan detail per merek dan opsi yang disarankan ada di tabel perbandingan merek pada Bab 13 "Alat Sistem Praktis".
  \item Platform penandaan utama: Penandaan nomor di Tiongkok terutama disediakan tiga penyedia data—Tencent Mobile Manager, 360 Mobile Guard, dan Dianhuabang (Tapir). Fitur "nama pemilik" dan "pemblokiran gangguan" dari operator (China Mobile/Unicom/Telecom) juga sering memakai data ini. Basis penandaan merek ponsel (Huawei, Xiaomi, OPPO, vivo) umumnya bekerja sama dengan penyedia di atas.
  \item Kalau salah ditandai: Jika nomor kerja Anda keliru ditandai sebagai "antar makanan", "kurir", atau "panggilan gangguan", Anda bisa mengajukan banding lewat kanal berikut: Dianhuabang: kunjungi platform banding nomor (haoma.tk.cn) dan lampirkan bukti identitas.
  \item 360: ajukan lewat fitur "Banding Nomor" di App 360 Mobile Guard.
  \item Tencent: ajukan lewat halaman "Banding Nomor" di mini program "Tencent Mobile Manager".
  \item Tiga operator besar: Pengguna China Mobile dapat mengirim SMS "KTFSR" ke 10086 untuk mengaktifkan fitur anti-gangguan; pengguna Telecom ikuti akun resmi "Tianyi Anti-Gangguan"; pengguna Unicom atur di "Unicom Mobile Manager".
\end{itemize}

Tingkat lanjut anti-gangguan: Aktifkan "pemblokiran panggilan frekuensi tinggi" dari operator (misalnya pembatasan panggilan frekuensi tinggi China Mobile) untuk otomatis memblokir panggilan keluar suara AI.

Kemampuan umum pemblokiran bawaan ponsel: Identifikasi panggilan / pemblokiran pintar: menandai dan memblokir panggilan gangguan, penipuan, dan promosi iklan berdasarkan basis data nomor.

Daftar hitam / daftar putih: blokir manual nomor gangguan yang diketahui; masukkan kontak serta nomor kurir/antar makanan ke daftar putih agar tidak salah blokir.

Filter kata kunci SMS: SMS yang memuat kata kunci seperti "menang hadiah", "pinjaman", "klik tautan" otomatis masuk kotak blokir.

iOS bisukan panggilan tak dikenal: panggilan yang tidak ada di kontak hanya dibisukan ke pesan suara, tidak mengganggu Anda; cocok untuk yang sering diganggu panggilan tak dikenal, tapi perlu waspada melewatkan panggilan penting dari nomor asing.

Terhubung dengan mode Jangan Ganggu: saat malam atau rapat, hanya izinkan panggilan dari kontak, lengkap dengan jadwal otomatis—kualitas tidur langsung membaik.

Cara memadukan perangkat lunak praktis: utamakan pemblokiran bawaan sistem, lalu tambahkan Pusat Anti-Penipuan Nasional dan layanan anti-gangguan operator; untuk pembersihan bisa pilih alat ringan seperti Qingzhuo, tetapi satu jenis perangkat lunak keamanan/pembersihan cukup satu set—lihat Bab 13.

\booksubsection{Apa itu VoLTE/VoNR dan cara mengaktifkannya:}

\begin{itemize}
  \item VoLTE (Voice over LTE): teknologi panggilan suara berbasis jaringan 4G. Saat menelepon, sinyal tidak turun ke 2G/3G lagi, jadi bisa menelepon sambil internetan, kualitas suara lebih jernih, dan tersambung lebih cepat.
  \item VoNR (Voice over NR): panggilan suara berbasis 5G, evolusi berikutnya dari VoLTE. Latensi lebih rendah, kualitas suara lebih baik, dan mendukung video call HD. Perlu ponsel 5G + paket 5G + cakupan 5G lokal agar bisa dipakai.
  \item Cara mengaktifkan: Di iOS, aktifkan "VoLTE" di "Pengaturan > Seluler > Opsi Data Seluler > Suara \& Data"; di Android, cari sakelar "VoLTE HD" di "Pengaturan > Jaringan Seluler > Data Seluler > Lanjutan". Operator biasanya sudah mengaktifkan secara bawaan; bila saat menelepon internet terputus dan tidak bisa online, kemungkinan VoLTE belum aktif.
\end{itemize}

\begin{tcbinfo}{Tips singkat}
Kalau saat menelepon internet tiba-tiba melambat atau putus, biasanya VoLTE/VoNR belum aktif dan jaringan turun ke 2G/3G. Setelah diaktifkan, Anda bisa “menelepon sambil nonton video”.

\end{tcbinfo}

\booksection{SMS dan iMessage/RCS}

SMS tradisional belum punah—masih menjadi kanal andalan untuk menerima kode verifikasi, notifikasi bank, dan pesan resmi.

\begin{itemize}
  \item iMessage: pesan antarperangkat Apple lewat internet, tampil sebagai gelembung biru. Mendukung foto/video berkualitas tinggi, tanda terima baca, sepenuhnya gratis, dan terenkripsi ujung ke ujung.
  \item RCS: protokol komunikasi generasi berikutnya yang didorong operator global dan Google. Fiturnya mirip iMessage atau WeChat, tapi berbasis nomor telepon.
  \item Interoperabilitas iMessage dan RCS: Sejak iOS 18, Apple mulai mendukung standar RCS. Artinya, iPhone dan ponsel Android bisa bertukar foto HD, tanda terima baca, dan indikator sedang mengetik lewat RCS—pengalaman "gelembung hijau" yang terputus pun berakhir. Syaratnya, kedua operator harus mengaktifkan RCS dan versi sistem ponsel cukup baru.
\end{itemize}

\booksubsection{Penerapan RCS di Tiongkok:}

\begin{itemize}
  \item China Mobile: paling awal menerapkan RCS; sebagian provinsi sudah mendukung layanan "Pesan 5G" untuk mengirim pesan media kaya dan kartu interaktif. Perlu ponsel yang mendukung dan diaktifkan di gerai operator.
  \item China Telecom: menyediakan layanan lewat "Telecom Pesan 5G"; sebagian paket sudah mendukung secara bawaan, pengguna lama perlu mengirim perintah SMS tertentu untuk mengaktifkan.
  \item China Unicom: menerapkan Pesan 5G di kota-kota utama, tetapi dukungan perangkat belum merata—banyak model Android masih perlu mengaktifkan sakelar "Pesan Lanjutan" atau "Pesan 5G" secara manual di App SMS.
  \item Pengalaman nyata: Adopsi RCS di Tiongkok lebih lambat dari perkiraan, terutama karena ekosistem WeChat yang kuat dan penyesuaian perangkat yang belum seragam. Komunikasi harian masih didominasi WeChat; RCS lebih banyak berperan sebagai "kanal resmi" (seperti notifikasi bank dan pesan pemerintah).
\end{itemize}

\begin{tcbwarn}{Garis merah keamanan kode verifikasi}
\begin{itemize}
  \item Jangan teruskan kode verifikasi: siapa pun yang meminta Anda "membacakan kode verifikasi" atau "meneruskan SMS" bisa jadi penipu—termasuk yang mengaku "layanan pelanggan", "kurir", atau "aparat penegak hukum". Kode verifikasi = dompet Anda.
  \item Waspadai SMS phishing: curigai tautan dalam SMS seperti "ETC kedaluwarsa", "kartu asuransi kesehatan dibekukan", "ganti rugi paket hilang". Lembaga resmi tidak akan meminta nomor kartu bank dan kata sandi lewat tautan SMS.
  \item Identifikasi nomor: SMS resmi bank biasanya berasal dari nomor pendek awalan 95 atau 1069 + tanda tangan perusahaan. "Notifikasi bank" yang dikirim dari nomor HP 11 digit biasanya hampir pasti penipuan.
  \item Cadangan dua faktor: jangan ganti nomor HP yang terikat ke akun penting sembarangan; jika terpaksa ganti nomor, pastikan melepas kaitan semua akun kunci terlebih dahulu, karena pemilik nomor baru bisa mengatur ulang kata sandi Anda lewat kode verifikasi SMS.
\end{itemize}

\end{tcbwarn}

\booksection{WeChat/QQ: Mengelola aplikasi sosial sejuta umat}

WeChat dan QQ sudah melampaui alat komunikasi; keduanya menjadi “super app” yang memadukan sosial, pembayaran, dan akses informasi.

\booksubsection{Pencadangan dan migrasi riwayat chat:}

\begin{enumerate}
  \item Masuk ke WeChat versi PC atau Mac di komputer.
  \item Klik menu di kiri bawah > [Migrasi \& Pencadangan] > [Cadangkan \& Pulihkan].
  \item Pilih "Cadangkan riwayat chat ke komputer", lalu pilih chat yang ingin dicadangkan di ponsel.
  \item Migrasi ke ponsel baru: di WeChat ponsel lama, buka [Saya] > [Pengaturan] > [Umum] > [Migrasi \& Pencadangan Riwayat Chat] > [Migrasi] > [Migrasi ke Ponsel/Perangkat Tablet], pindai kode QR dan selesai.
\end{enumerate}

\booksubsection{Cara mendalam membersihkan ruang penyimpanan:}

\begin{itemize}
  \item Jalur umum: WeChat [Saya] > [Pengaturan] > [Umum] > [Ruang Penyimpanan]—bisa membersihkan cache dan riwayat chat secara detail berdasarkan waktu serta jenis file.
  \item Pembersihan per chat: buka sebuah chat > "…" di kanan atas > [Riwayat Chat] > [Kelola Ruang Penyimpanan]—bisa menghapus foto, video, dan file di chat tertentu sambil mempertahankan teks.
  \item Bersihkan foto/video asli: foto dan video HD asli di grup paling memakan ruang. Di [Kelola Ruang Penyimpanan], centang "Foto asli" untuk menghapus massal.
  \item Matikan unduhan otomatis: [Pengaturan] > [Umum] > [Foto, Video, File, dan Panggilan], matikan "Unduhan otomatis" agar tidak diam-diam meng-cache banyak file grup di WiFi.
  \item Cache mini program: di [Temukan] > [Mini Program], tekan lama mini program yang jarang dipakai lalu hapus untuk membebaskan cache lokalnya.
  \item Pencadangan terkompresi riwayat chat: saat mencadangkan lewat WeChat komputer, centang "Kompres" untuk memampatkan riwayat chat ber-GB menjadi 30\%–50\% dari ukuran aslinya.
\end{itemize}

\booksubsection{Tips mengelola grup:}

\begin{itemize}
  \item Pengumuman grup: pemilik/admin grup memublikasikannya di "…" kanan atas chat grup > [Pengumuman Grup]. Semua anggota menerima notifikasi kuat—cocok untuk pemberitahuan penting. Otomatis muncul saat anggota baru bergabung.
  \item Tugas grup: pemilik grup dapat mengatur pesan yang disematkan sebagai "Tugas Grup"; semua anggota melihat penanda merah di atas, dan hilang otomatis setelah diklik selesai. Cocok untuk pendaftaran acara dan pengumpulan tugas.
  \item Rantai Grup (Jielong): ketik "\#jielong" di kotak input untuk memulai rantai, cocok untuk pesan makan bersama atau pengumpulan daftar pendaftaran. Hasil rantai otomatis dirangkum di bagian atas grup, sehingga chat tidak banjir pesan.
  \item Siaran langsung grup: " + " di kanan atas chat grup > [Siaran Grup]—bisa memulai live yang hanya bisa ditonton anggota grup; cocok untuk pelatihan internal, kumpul keluarga, atau uji coba produk internal.
  \item Lipat chat grup: di daftar grup, geser nama grup ke kiri > [Lipat]—grup yang kurang penting bisa dirapikan ke "Grup Terlipat", dan hanya memberi notifikasi saat ada yang @ Anda.
  \item Hanya pemilik/admin yang boleh @ semua: saat grup berisi 500 orang ke atas, anggota biasa tidak bisa @ semua, agar tidak saling mengganggu.
\end{itemize}

\booksubsection{Pengelolaan izin Moments secara detail:}

\begin{itemize}
  \item Visibilitas kelompok: saat membagikan, klik "Siapa yang boleh melihat", pilih "Terlihat sebagian" atau "Sembunyikan dari"—berdasarkan label kelompok yang sudah diatur (misalnya "Rekan kerja", "Keluarga", "Klien").
  \item Rentang teman yang boleh melihat: [Saya] > [Pengaturan] > [Izin Teman] > [Moments] > [Rentang teman yang boleh melihat Moments]—pilih "3 hari terakhir / 1 bulan / 6 bulan / Semua". Disarankan pilih "6 bulan" untuk menyeimbangkan privasi dan sosial.
  \item Sembunyikan satu unggahan Moments: tekan lama unggahan di Moments > [Jangan lihat Moments-nya] atau [Berhenti ikuti]—pihak lain tidak akan tahu.
  \item Kelola iklan Moments: [Saya] > [Pengaturan] > [Tentang WeChat] > [Informasi Pribadi \& Izin] > [Kelola Iklan Personal]—bisa mematikan rekomendasi personal berbasis perilaku.
  \item Sepuluh unggahan untuk orang asing: [Izin Teman] > [Moments] > [Izinkan orang asing melihat sepuluh unggahan Moments]—matikan agar orang asing tidak bisa menelusuri unggahan lama Anda.
\end{itemize}

\booksection{Perbandingan video call}

Setiap alat video call punya keunggulan masing-masing. Memilih alat yang sesuai dengan situasi membuat komunikasi lebih lancar.

\begin{itemize}
  \item Video WeChat: jangkauan sejuta umat, cocok untuk semua usia. Ideal untuk kontak harian keluarga dan teman, komunikasi cepat. Maksimal 20 orang dalam video (perlu mini program "Grup Video"); kekurangannya, kualitas gambar turun drastis saat jaringan buruk.
  \item FaceTime: khusus ekosistem Apple, kualitas gambar sangat baik, terenkripsi ujung ke ujung. Sejak iOS 15 mendukung undangan "tautan" untuk pengguna Android/Web, tetapi host harus perangkat Apple. Cocok untuk komunikasi keluarga lintas perangkat.
  \item Zoom: pilihan utama untuk rapat internasional, mendukung rapat gratis hingga 100 orang selama 40 menit. Fitur berbagi layar, latar virtual, dan ruang diskusi kelompok lengkap—cocok untuk kolaborasi lintas negara/jarak jauh.
  \item Tencent Meeting: pilihan utama untuk kerja kantoran di Tiongkok, terintegrasi erat dengan WeChat, masuk cukup dengan pindai QR. Mendukung rapat 300 orang (versi berbayar); perekaman cloud, takarir langsung, dan notulen rapatnya kuat—cocok untuk pelatihan internal dan rapat besar.
  \item Saran pemakaian: kontak harian 1 lawan 1 pilih WeChat; antarpengguna Apple pilih FaceTime; rapat bisnis resmi pilih Zoom atau Tencent Meeting (di Tiongkok Tencent, lintas negara Zoom).
\end{itemize}

\booksubsection{Kebutuhan bandwidth dan perkiraan konsumsi kuota:}

\begin{itemize}
  \item Video call WeChat: sekitar 500 Kbps–1 Mbps; kualitas standar memakan sekitar 3–5 MB per menit, HD sekitar 8–12 MB.
  \item FaceTime: SD sekitar 1 Mbps, 720p sekitar 2 Mbps, 1080p butuh 4 Mbps ke atas. Kuota sekitar 5–15 MB per menit.
  \item Zoom/Tencent Meeting: 720p sekitar 1,5–3 Mbps, 1080p butuh 3–5 Mbps. Kuota sekitar 300–800 MB per jam; host rapat multi-peserta memakan lebih banyak.
  \item Saran: untuk video call lama, usahakan pakai WiFi; di jaringan seluler, turunkan kualitas gambar secara manual (matikan "Video HD" di pengaturan App) untuk menghemat kuota.
\end{itemize}

\booksubsection{Tips mengoptimalkan kualitas video call:}

\begin{itemize}
  \item Pencahayaan: punggung menghadap jendela, wajah menghadap sumber cahaya (cahaya alami terbaik); hindari cahaya dari belakang yang membuat wajah gelap.
  \item Sudut: letakkan kamera ponsel sedikit lebih tinggi dari mata agar dagu ganda tidak terlihat; saat panggilan lanskap, pegang ponsel dengan dua tangan lebih stabil.
  \item Audio: gunakan headset untuk menghindari gema dan dengung; di lingkungan tenang, jangan aktifkan mode "Speaker".
  \item Jaringan: utamakan WiFi; saat pakai 4G/5G, dekati jendela agar sinyal tidak melemah; di rapat multi-peserta, bisa matikan video orang lain dan tampilkan hanya pembicara.
  \item Latar buram: Tencent Meeting dan Zoom mendukung buram AI untuk menyamarkan latar yang berantakan. Kenakan pakaian polos agar tidak mirip warna latar dan menimbulkan efek "kunci hijau" yang janggal.
\end{itemize}

\begin{tcbinfo}{Tips kecil}
Saat jaringan tersendat, matikan video dan pertahankan suara—lebih sopan daripada langsung menutup panggilan. Sebelum rapat penting, uji bandwidth lebih dulu dengan fitur “Deteksi Jaringan” bawaan Tencent Meeting.

\end{tcbinfo}

\booksection{Email: konfigurasi dan pemakaian di ponsel}

\begin{itemize}
  \item Menambahkan akun email: di iOS, buka Pengaturan > Mail > Akun > Tambah Akun; di Android biasanya tersedia App "Email".
  \item Tips pemakaian: atur notifikasi berbeda untuk setiap akun email; manfaatkan gestur geser untuk cepat mengarsipkan, menghapus, atau menandai email.
\end{itemize}

\booksubsection{Perbedaan protokol IMAP vs POP3 vs Exchange:}

\begin{itemize}
  \item POP3: setelah mengunduh email dari server, salinan asli dihapus secara bawaan dan hanya tersimpan di perangkat. Cocok untuk satu perangkat dan kebutuhan baca offline. Kekurangannya, setelah ganti perangkat Anda tidak melihat email lama—sudah tidak direkomendasikan.
  \item IMAP: email selalu tersimpan di server, di perangkat hanya cache. Sinkronisasi multi-perangkat (ponsel/komputer/tablet) konsisten—sudah dibaca di satu perangkat, dianggap sudah dibaca di semua. Protokol utama yang direkomendasikan saat ini.
  \item Exchange (termasuk ActiveSync): protokol Microsoft yang, selain email, juga menyinkronkan kalender, kontak, dan tugas. Pilihan utama untuk email perusahaan. Didukung Outlook, NetEase Enterprise Mail, dan Tencent Enterprise Mail.
  \item Saran pemilihan: email pribadi pakai IMAP; email kantor utamakan Exchange, lalu IMAP; sebisa mungkin jangan pakai POP3 kecuali ada kebutuhan arsip khusus.
\end{itemize}

\booksubsection{Cara mengonfigurasi email perusahaan:}

\begin{enumerate}
  \item iOS: [Pengaturan] > [Mail] > [Akun] > [Tambah Akun] > pilih "Microsoft Exchange" atau "Lainnya".
  \item Android: di App "Email" bawaan atau App Outlook, pilih "Tambah akun", masukkan alamat email lengkap dan kata sandi.
  \item Pengisian server (contoh Exchange): alamat server biasanya "mail.namadomainperusahaan.com" atau "outlook.office365.com"; isian domain biasanya singkatan nama perusahaan.
  \item Jika muncul pesan "tidak dapat memverifikasi server", isi manual: server masuk (IMAP) port 993, enkripsi SSL; server keluar (SMTP) port 465 atau 587, enkripsi TLS.
  \item Setelah masuk, disarankan aktifkan "Push" di [Pengaturan Akun] agar email diterima real-time, dan matikan "Pratinjau Email" untuk melindungi rahasia dagang.
\end{enumerate}

\begin{tcbinfo}{Catatan kecil}
Email perusahaan biasanya mengharuskan admin mengaktifkan “kata sandi otorisasi klien” atau “kode otorisasi IMAP/SMTP”—tidak bisa langsung memakai kata sandi login web. Jika gagal masuk, buat kode otorisasi di pengaturan email versi web, lalu masukkan ke ponsel.

\end{tcbinfo}

\begin{tcbpractice}{Praktik: cadangkan riwayat chat WeChat ke komputer}
\begin{enumerate}
  \item Mulai pencadangan di komputer: buka WeChat versi PC > menu kiri bawah > [Migrasi \& Pencadangan] > [Cadangkan \& Pulihkan].
  \item Pilih cadangkan ke komputer: klik "Cadangkan riwayat chat ke komputer".
  \item Pilih konten cadangan di ponsel, lalu tunggu hingga selesai.
\end{enumerate}

\end{tcbpractice}
